\section{Adjusting Exponents in Frozen Models}
\label{sec:mature-adjustments}
\label{sec:exp-frozen-support}
\label{sec:exp-zero-training-system}
\label{sec:coordinate}

The crossing experiments in \S\ref{sec:coadaptation-finding} show that trained weights are co-adapted to their specific exponent basis, motivating adjustments defined relative to native exponents.

\subsection{A common coordinate for exponent adjustments}

For native frequencies $\omega_k^N$ and installed frequencies $\omega'_k$, define
\begin{equation}
d_k=\log\frac{\omega_k^N}{\omega'_k},\qquad
\omega'_k=\omega_k^N e^{-d_k},\qquad
\phi'_k-\phi_k^N=\frac{d_k}{\log b}.
\label{eq:movement-allocation}
\end{equation}
\begin{equation}
d_k=
\begin{cases}
-\log(1-w_k+w_k/s), & \text{frequency blend},\\
m_k\log s, & \text{log-frequency shift},\\
\sum_{j<k}\log\lambda_j, & \text{mixed-radix conversion}.
\end{cases}
\label{eq:adjustment-families}
\end{equation}
Static tables are installed before prefill and retained through decoding; amplitude $g$ scales both cosine and sine tables ($g^2$ on attention logits).

\subsection{Fixed-range and single-table comparisons}

\paragraph{Frozen fixed-support control.}
On frozen OLMo-$2$-$1$B and Qwen-$2.5$-$1.5$B, we hold endpoints, $64$ pairs, attention amplitude, inputs, and decoder fixed. Geometric, coarse label-free ramp, and residual-guided allocations give the results in Table~\ref{tab:frozen-exponent-main}. Structured profiles produce large gains over uniform exponents (OLMo $16$K RULER: $0.56\%$ to $61.04\%$), with coarse controls capturing the bulk of the effect (Appendix~\ref{sec:frozen-fixed-support}).

\begin{table}[t]
\centering
\small
\caption{\textbf{Frozen-model RULER scores (\%) at matched exponent range.}
OLMo uses a held-out nine-task confirmation at $16$K; Qwen uses a four-task
development panel at $64$K ($20$ rows per task).}
\label{tab:frozen-exponent-main}
\begin{tabular}{@{}lrr@{}}
\toprule
Exponent allocation & OLMo-$1$B & Qwen-$1.5$B \\
\midrule
Geometric & 0.56 & 57.75 \\
Coarse ramp & \textbf{61.04} & 64.00 \\
Residual-guided & 60.47 & \textbf{66.50} \\
\bottomrule
\end{tabular}
\end{table}

\paragraph{One static configuration across lengths.}
A normalized-index profile transfers a reference displacement curve onto the target grid ($\omega'_k=\omega_k^N s^{-m_k}$). On Qwen-$2.5$-$0.5$B, the same factor-two table is evaluated at $32$K and $64$K across full RULER-13 ($20$ rows per task). With reported fixed amplitudes, Index achieves $55.92\%$ and $51.46\%$, compared with $55.94\%$ and $45.37\%$ for YaRN (Table~\ref{tab:static-index-main}). The $64$K improvement is $6.09$ points ($95\%$ bootstrap interval $[2.76, 9.58]$), while $32$K scores are near-equal. The gain is attributed to the complete static configuration (table $\times$ amplitude; Appendix~\ref{sec:index-adjustment}).

\begin{table}[t]
\centering
\small
\caption{\textbf{Static extension across native and extended windows.}
Qwen-$0.5$B RULER-13 macro (\%); $g$ is fixed amplitude.}
\label{tab:static-index-main}
\begin{tabular}{@{}lrrr@{}}
\toprule
Installed table & $g$ & $32$K & $64$K \\
\midrule
Native & 1.0000 & 54.78 & 22.05 \\
YaRN-$2$ & 1.0693 & 55.94 & 45.37 \\
Normalized-index-$2$ & 1.0513 & 55.92 & \textbf{51.46} \\
\bottomrule
\end{tabular}
\end{table}

\paragraph{Finite-grid placement.}
Comparing normalized-index interpolation against direct-gap evaluation on Gemma-$1.1$-$2$B ($128$ pairs; \citealt{gemmateam2024gemma}) yields $79.00\%$ versus $72.81\%$ at $16$K ($+6.19$ points, $[2.81, 9.63]$), identifying placement as an important implementation choice (Appendix~\ref{sec:coordinate-confirmation}).

\subsection{Redistributing the intermediate-band shift}
\label{sec:bm-main}

MrRoPE-Pro progressively distributes the cumulative exponent shift across an intermediate band $k\in[l,h]$ with $N=h-l$ and $q=\operatorname{clip}(k-l,0,N)$:
\begin{equation}
m_q^{\rm Pro}=\frac{q(q+1)}{N(N+1)},
\qquad
m_q^{\rm BM}=\frac{q(q+1)(3N+2-2q)}{N(N+1)(N+2)}.
\label{eq:bm-formula}
\end{equation}
Boundary matching (BM) minimizes quadratic difference energy, smoothing radix transitions toward unmodified outer bands (Appendix~\ref{sec:bm-construction}).

\begin{figure}[t]
  \centering
  \includegraphics[width=\linewidth]{fig_bm_natural_qa.pdf}
  \caption{\textbf{Intermediate-band redistribution improves OLMo natural QA.}
  Complete-output token F1 on untruncated long inputs ($s=4$). Sample counts shown by task; final row averages task means.}
  \label{fig:bm-natural-qa}
\end{figure}

In the OLMo-$1$B confirmation, BM reaches $51.32\%$ at $16$K compared with $2.78\%$ for MrRoPE-Pro. Crucially, across $631$ untruncated LongBench prompts exceeding $4$K tokens across five tasks, BM raises task-equal complete-output token F1 from $21.62\%$ to $25.44\%$ (Fig.~\ref{fig:bm-natural-qa}; $95\%$ bootstrap interval $[1.32, 6.29]$, all task differences positive).

\begin{table}[t]
\centering
\small
\setlength{\tabcolsep}{4pt}
\caption{\textbf{Matched six-task RULER comparisons (\%).}
Each length column averages six task scores ($s=4$; short/long prompt counts per arm).}
\label{tab:bm-models}
\begin{tabular}{@{}lclrrrr@{}}
\toprule
& & & \multicolumn{2}{c}{Short} & \multicolumn{2}{c}{Long}\\
Model & Lengths & Rows & MrPro & BM & MrPro & BM \\
\midrule
OLMo-$1$B & $4/16$K & $24/48$ & 37.85 & 81.81 & 2.78 & 51.32 \\
Qwen-$3$B & $32/128$K & $12/24$ & 87.22 & 91.67 & 78.13 & 70.83 \\
Qwen-$7$B & $32/128$K & $6/12$ & 83.33 & 80.00 & 84.44 & 71.11 \\
\bottomrule
\end{tabular}
\end{table}

However, preferences differ across models (Table~\ref{tab:bm-models}): while OLMo favors BM, Qwen-$3$B and $7$B favor MrRoPE-Pro at $128$K ($78.13\%$ vs $70.83\%$ and $84.44\%$ vs $71.11\%$). With endpoints and total scale fixed, intermediate profile preferences are checkpoint- and length-dependent (Appendix~\ref{sec:bm-model-comparison}).
